\subsection{Background: The Agent Capability Governance Gap}\label{subsec:background}

Large language model (LLM) agents are proliferating across ecosystems. Each agent exposes a set of \emph{capabilities}: tools, APIs, knowledge domains, and interaction protocols. A capability is a claim---``I can retrieve weather data,'' ``I can execute SQL,'' ``I can reason about financial risk''---that other agents and human operators must evaluate before delegation. But who records what a capability does, how it evolves, and whether it is trustworthy? Today, capability claims are scattered across README files, API documentation, and prompt templates. They are revised without audit trails, deprecated without ceremony, and forked without attribution. A lightweight, verifiable, lifecycle-aware registry of agent capabilities does not yet exist.

\paragraph{Existing governance layers.}
Three recent lines of work address adjacent aspects of agent governance, but none fills the capability-registry gap.

\textbf{KYA} (\emph{Know Your Agent}, arXiv:2605.25376) is a framework-agnostic trust layer for agent systems that provides a permissions model with HMAC-chained integrity \cite{kya2025}. It achieves approximately 1,800 operations per second and is available as the \texttt{veldt-kya} package on PyPI. KYA governs \emph{who} can invoke \emph{what}, but it does not govern the \emph{lifecycle} of a capability itself: there is no \texttt{VALIDATE} event that records an independent audit of a tool's behavior, no \texttt{DEPRECATE} event that marks a retired protocol, and no \texttt{FORK} event that traces a community derivative. KYA secures the permission boundary; it does not register the capability's provenance.

\textbf{AgentSafe} (arXiv:2512.03180) is a unified governance framework spanning design-time, runtime, and audit phases \cite{agentsafe2025}. It provides architecture-level governance---data-flow validation, policy enforcement, and safety guardrails---but it is heavyweight and requires integration with a service mesh or runtime orchestrator. It does not offer a lightweight, Markdown-native format that a single agent can author, cryptographically link, and commit to version control without enterprise infrastructure.

\textbf{LLM-native ontology engineering} (Talukdar et al., arXiv:2604.23090; LLM4ACOE) uses multi-agent LLM pipelines to generate OWL ontologies from natural-language requirements \cite{talukdar2025}. These systems produce static taxonomies and relation schemas, but they do not govern the \emph{runtime lifecycle} of the concepts they create: there is no event-sourced record of validation, deprecation, or fork, and no tamper-evident chain linking an agent's assertion to its subsequent revision or retirement. The ontology is produced, not governed.

Two additional bodies of work are complementary but distinct. \textbf{``From Agent Traces to Trust''} (arXiv:2606.04990) surveys provenance frameworks for agent systems and explicitly notes that ``a unified trace schema is still needed'' \cite{agenttraces2025}. ADL Lite is that schema: an event-first, cryptographically linked trace format in which every capability claim, validation, deprecation, and fork is a first-class auditable event. \textbf{SafeAgent} (arXiv:2604.17562) is a safety protocol for LLM agents that defines red-team evaluation procedures and harm-classification taxonomies \cite{safeagent2025}. It governs \emph{safety assessment}, not \emph{capability registration}; its concern is whether an agent is dangerous, not whether its advertised capabilities are authentic, current, and traceable.

\paragraph{Defining the gap.}
No existing system simultaneously satisfies four properties that define the capability-registry design space: (a)~document-native authoring in plain Markdown, (b)~built-in tamper detection through cryptographic hash chaining, (c)~lifecycle state machines with declarative preconditions, and (d)~pip-installable deployment requiring no enterprise infrastructure. KYA satisfies (b) and (d) but lacks (a) and (c). AgentSafe partially satisfies (c) but lacks (a), (b), and (d). LLM-native OE tools satisfy (a) in a limited sense (natural-language input) but lack (b), (c), and (d). The intersection of all four properties remains unoccupied.

\paragraph{Defining ``operational ontology.''}
We use the term \emph{operational ontology} in a specific sense that requires explicit definition, because the term carries distinct meanings in adjacent literatures. In the OBO Foundry ecosystem, an ``operational ontology'' typically means an application ontology deployed in a running bioinformatics pipeline (as opposed to a reference ontology) \cite{obo_foundry}. In programming-language theory, ``operational semantics'' defines the meaning of a program through its execution on an abstract machine. In knowledge-engineering practice, ``executable ontologies'' (e.g., UFO-B executable specifications \cite{guizzardi2024ufob}) are ontologies equipped with first-order logic axioms that can be evaluated by a reasoner. ADL Lite's usage differs from all three: we mean an ontology in which the \emph{operational semantics of lifecycle transitions}---preconditions, state derivation, and consensus rules---are co-located with the \emph{data structure} (EventChain) in a single lightweight package. In the present framing, the ``ontology'' is the set of capability concepts an agent ecosystem manages; the ``operational'' aspect is the fact that the same YAML-declared rules that define what can happen to a capability are enforced at append time by the ActionExecutor, without requiring an external workflow engine, OWL reasoner, or triple store. This co-location of governance semantics and data structure is the defining characteristic of ADL Lite's operational ontology.

\paragraph{Ontological novelty of the event-first stance.}
We emphasize that ADL Lite's contribution is not the philosophical thesis that ``events are fundamental''---a position already well-established by Wittgenstein's \emph{Tractatus} (\S1.1), BFO's \emph{occurrent}, DOLCE's \emph{perdurant}, and UFO-B's \emph{event} \cite{wittgenstein_tractatus,bfo,dolce,ufo}. Rather, ADL Lite's contribution is the \emph{operationalization} of this stance: the translation of an event-first philosophical commitment into a deployable, Markdown-native, cryptographically enhanced capability-lifecycle mechanism. In foundational ontologies, events and continuants are typically treated as co-equal top-level categories; ADL Lite treats events as the \emph{sole architectural primitive} from which all lifecycle state---status, confidence, validator identity---is derived through deterministic functions. This is an engineering reduction, not a metaphysical elimination: capabilities are not denied ontological standing, but their \emph{computable} properties (status, confidence) are made dependent on event history. This operationalization distinguishes ADL Lite from descriptive event ontologies (CIDOC CRM, SEM/LODE) that record what happened without enforcing what can happen. It also distinguishes ADL Lite from cryptographic provenance systems (nanopublications, PROV-O) that provide tamper detection without lifecycle state machines.

\paragraph{Why existing families are insufficient.}
We demonstrate systematic awareness by showing that representative systems from five distinct families fail to cover the full intersection of properties (a)--(d):

\textbf{Descriptive event ontologies (CIDOC CRM, SEO).} CIDOC CRM provides event-centric descriptive modeling and is widely deployed in cultural heritage \cite{cidoc_crm}. The Subject-Event Ontology (SEO) offers event-first patterns for web-scale data \cite{boldachev2025seo}. Both fail to satisfy (b) and (c): they record what happened without cryptographic integrity guarantees and without declarative preconditions for lifecycle state transitions. Moreover, production deployments typically require reasoning-enabled RDF/OWL toolchains, triple stores, and SPARQL endpoints, violating (d).

\textbf{Cryptographic provenance systems (PROV-O, nanopublications with Trusty URIs).} PROV-O provides a W3C-standard provenance vocabulary \cite{w3c_prov_o}; nanopublications with Trusty URIs achieve cryptographic integrity for atomic claims through hash-based URIs \cite{nanopublications,trusty_uris}. However, PROV-O lacks (b) natively: it describes provenance without built-in tamper detection, and it requires RDF tooling that conflicts with (a) and (d). Nanopublications satisfy (b) for atomic claims but lack (a) because they require RDF serialization and lack (c) because they provide no built-in lifecycle governance or state-machine preconditions.

\textbf{Markdown-native knowledge tools (Foam, Obsidian).} These tools satisfy (a) and (d) by design: they operate on plain Markdown files with minimal installation overhead \cite{foam2020,obsidian2024}. However, they provide neither (b) nor (c): they lack cryptographic hash chaining and provide no lifecycle state machines or precondition enforcement.

\textbf{Summary.} Each family lacks at least one property, and no family satisfies all four. ADL Lite occupies this intersection by integrating these properties into a coherent operational ontology, not by inventing any individual component.

\subsection{Research Gap: Capability Lifecycle Governance as Ontological Analysis}\label{subsec:research-gap}

The specific research gap addressed in this paper can be formulated as follows: \emph{How can an event-first capability registry achieve lifecycle traceability and multi-agent governance without requiring object-first mutable state, expert curation, or heavyweight tooling?} This gap has both an engineering dimension and an ontological dimension.

The \textbf{engineering dimension} concerns deployment friction. Existing operational ontology platforms require specialized expertise, dedicated infrastructure, and manual curation workflows that do not scale with the velocity of agent-generated capability artifacts \cite{hogan2021knowledge,guha2016schemaorg}. LLM-based agents cannot participate in capability registry construction because they lack programmatic access to the proprietary interfaces through which ontologies are constructed.

The \textbf{ontological dimension} concerns the status--history relationship. In an object-first registry, a capability's status is a property of the capability record. In an event-first registry, a capability's status is a derivative of its history. The transition from object-first to event-first is not merely a software design pattern (though it resembles event sourcing \cite{event_sourcing,cqrs}); it is an ontological reorientation that changes the fundamental categories of the registry. While event-centric modeling has been extensively analyzed in the applied ontology literature (CIDOC CRM, SEM, SEO, UFO-B), the \emph{operationalization} of event-first semantics---combining cryptographic integrity, deterministic lifecycle derivation, and declarative preconditions in a single lightweight system---has not been systematically studied.

This paper addresses both dimensions by presenting ADL Lite, an event-first operational ontology in which capabilities are modeled as append-only, cryptographically linked EventChains, and all lifecycle state is derived exclusively from event history. The paper's primary contribution is an \textbf{ontological analysis} that situates ADL Lite within the landscape of foundational ontologies (BFO, DOLCE, UFO), formalizes its categories and identity conditions, and proves key properties of its derivation semantics.

\subsection{Contributions}\label{subsec:contributions}

The contributions of this paper are organized around three axes: ontological analysis, formal semantics, and empirical validation.

\textbf{Contribution 1: Ontological analysis of event-first capability governance.} We analyze ADL Lite's core categories through the lens of three influential foundational ontologies:
\begin{itemize}[leftmargin=2em]
    \item We show that ADL Lite's \texttt{Event} corresponds to the BFO category of \emph{occurrent} and to the DOLCE category of \emph{perdurant}.
    \item We argue that ADL Lite's \texttt{Concept} (here, a capability definition) is best understood as a BFO \emph{generically dependent continuant} or as a DOLCE \emph{non-physical object}.
    \item We demonstrate that ADL Lite's L3 relations instantiate UFO \emph{relators} with event-based lifecycles.
    \item We formalize identity conditions for events, capabilities, and chains, and analyze three forms of ontological dependence (historical, generic, mutual).
\end{itemize}

\textbf{Contribution 2: Formal derivation semantics with proved properties.} We provide a complete formal foundation comprising:
\begin{itemize}[leftmargin=2em]
    \item An event alphabet $\Sigma = \Sigma_{\text{life}} \cup \Sigma_{\text{comm}}$ partitioned into lifecycle and communication events.
    \item A well-formedness predicate $\text{WF}(C)$ over EventChains.
    \item Status derivation $\delta(C)$ and confidence aggregation $\gamma(C)$ as deterministic functions.
    \item Fork and confluence semantics.
    \item Seven proved properties: determinism (Theorem~1), confluence under fork (Theorem~2), transition monotonicity (Theorem~3), confidence boundedness (Theorem~4), confidence monotonicity under independent validation (Theorem~5), status--confidence consistency (Theorem~6), and CRDT convergence (Theorem~7).
    \item An analysis of expressive power in relation to Event Calculus and Description Logic.
\end{itemize}

\textbf{Contribution 3: Empirical validation of formal semantics and ontological consistency.} We evaluate ADL Lite on 201 EventChains derived from the IBM Anti-Money Laundering (AML) HI-Small dataset (9,300 events), empirically validating the formal semantics against seven proved theorems from Contribution~2. The evaluation confirmed:
\begin{itemize}[leftmargin=2em]
    \item Zero integrity failures across all chains, including the longest chain (300 events).
    \item 100\% accuracy on 2,204 exhaustive status derivation test cases.
    \item Perfect precision and recall (P = R = F1 = 1.0) on 19 precondition enforcement test cases.
    \item Linear scalability: 20{,}955 events per second on modest hardware.
\end{itemize}
These results are not a domain-level evaluation of AML effectiveness; they validate that the event-first ontological architecture---including the deterministic derivation semantics, hash-linked integrity guarantees, and seven formally proved properties---is computationally feasible, cryptographically sound, and empirically correct. Domain-level evaluation (E5, E6) is ongoing and will be reported in a future publication; the present paper is positioned as a \emph{framework paper} that establishes the ontological and formal foundations.

\subsection{Roadmap and Phase Definitions}\label{subsec:roadmap}

ADL Lite is developed in three phases with distinct objectives, threat models, and completeness criteria:

\begin{itemize}[leftmargin=2em]
    \item \textbf{Phase~1 (v0.1--v0.2, current):} Collaborative-audit trust model. Agents are assumed non-Byzantine; actor identity is self-declared string identifiers; fork resolution uses last-write-wins. This phase establishes the core EventChain data structure, deterministic derivation semantics, and precondition system. All empirical validation in this paper (E1--E3, E4, E6) falls within Phase~1.
    \item \textbf{Phase~2 (v0.3--v0.4, planned):} Structured agent communication. L2 Markdown body gains template-governed sections ([Observation], [Reasoning], [Conclusion]) with SHACL validation. Controlled forgetting (\textsc{ARCHIVE} events, embedding-based deduplication) addresses append-only accumulation. Confidence calibration via MARGIN EWMA reduces overconfidence.
    \item \textbf{Phase~3 (v0.5+, planned):} Adversarial robustness. W3C Linked Data Proofs (Ed25519 signatures + DIDs) replace string identifiers. Merkle-tree batch verification reduces audit complexity to $O(\log n)$. Optional BFT transport layer for Byzantine fault tolerance. Domain-level evaluation (E5, E6) with AML experts is targeted for Phase~3.
\end{itemize}

The governance claims in this paper are scoped to Phase~1: collaborative audit among non-Byzantine agents. Byzantine robustness, cryptographic authentication, and adversarial resilience are planned for Phase~3 and scoped to future work. The present paper reports Phase~1 results; Phase~2 and Phase~3 are in scope for future work (Section~\ref{subsec:future-work}).

\subsection{Paper Organization}\label{subsec:organization}

Section~\ref{sec:related-work} positions ADL Lite within the emerging agent governance landscape (KYA, AgentSafe, Agent Traces to Trust, and SafeAgent) and surveys related work in foundational ontologies, ontology engineering, and event-centric modeling. Section~\ref{sec:ontological-analysis} presents the core ontological analysis: categories, identity conditions, ontological dependence, and alignment with BFO, DOLCE, and UFO. Section~\ref{sec:architecture} describes the ADL Lite architecture, formal semantics, and comparison with Event Calculus and Description Logic. Section~\ref{sec:empirical-validation} reports empirical validation of architectural correctness. Section~\ref{sec:discussion} discusses implications, limitations, and future work. Section~\ref{sec:conclusion} concludes.
